% Single-contract dynamics: sample construction and estimators.
% Code: src/albert/research/pm_vol.py (steps prep, fit, vi, fc, hf_extract, hf, drift_sim, drift_tail).

\subsubsection{Sample, clock time and look-ahead guards}\label{sec:dyn-sample}

\paragraph{Hourly panel.}
We use the hourly candles of \cref{sec:data} for the settled markets of the seven non-financial
categories, February to October 2026. Kalshi omits a candle when nothing changed, so each market
is placed on a regular hourly grid from its first candle to its last candle at or before
\texttt{close\_time}, with the closing yes bid and ask forward-filled. Candles ending after the
close are dropped. This gives 4{,}828{,}986 market-hours on 47{,}505 markets.
% source: data/research/pm_vol/prep.log
Sports rows are split into \emph{pre-game} (the step ends before the scheduled start, parsed from the
ticker or set to expected expiration minus 3\,h) and \emph{in-game}.

A step is $(t,t+h]$ with $h=1$\,h. The decision quote must have a spread of at most 4\cents{} and a
mid in $[0.03,0.97]$; the end quote must be two-sided with a spread of at most 10\cents. This
measurement filter removes 0.4--5.2\% of otherwise valid steps.
% source: results_cat.json:categories.*.dropped_end_quote / base_steps
Mids are clipped to $[0.005,0.995]$ before taking $\Zt=\Phi^{-1}(\pr)$.

\paragraph{Time to the scheduled end.}
The model needs $\tauexp$, the time until information stops arriving (\cref{sec:model}).
\texttt{close\_time} cannot be used: for markets that close early it is look-ahead, since an early
close usually reveals the outcome. \texttt{expected\_expiration\_time} is the settlement schedule
and can lag the information by days. We therefore set $T=\texttt{expected\_expiration\_time}-o_s$
per series. The offset $o_s$ is the modal (expiration $-$ close) among the series' markets that
closed before 1 October 2026 if at least half share it, and otherwise their median over NO-settled
markets, which run to the scheduled end. Series without such markets have no schedule and are
excluded; this removes all October Elections series, so Elections has no holdout. Sports use
$T=\texttt{expected\_expiration\_time}$, about the start plus 3\,h.

On NO-settled markets the median of $\texttt{close\_time}-T$ is 0 in every non-sports category in
the \estim{}, and in the holdout everywhere except Mentions ($+10.4$\,h).
% source: report research-pm-vol-20261109T215533Z (scheduled-end table)
Steps ending less than 30 minutes before $T$ are dropped. The \estim{} (before 1 October 2026)
contains steps that end before 1 October and the holdout (October 2026) contains steps that start on
or after it. All parameters and tuning choices use the \estim{} only, and $t$-statistics are
clustered by event (\cref{sec:method}).

\subsubsection{Clock, innovations and the variance identity}\label{sec:dyn-clock}

Let $L_t=\int_t^{t+h}\irate_s\,\dd s=\ln(\Vrem_t/\Vrem_{t+h})$ be a step's clock increment, so that
M0 gives $\Zt_{t+h}\mid\Zt_t\sim\N(\Zt_t e^{L_t/2},\,e^{L_t}-1)$ (\cref{sec:model}). Treating the
observed mid as the fair price plus noise of variance $k s_t^2$ ($s_t$ the spread), we maximize
\begin{equation}\label{eq:dyn-nll}
\ell=\sum\log\varphi_{\N}\!\Big(\tilde\Zt_{t+h};\,\tilde\Zt_t e^{L_t/2},\;e^{L_t}-1+k\,s_t^2/\varphi(\tilde\Zt_t)^2\Big)
\end{equation}
over the clock and $k$. Tildes denote observed mids.

The candidate clocks are uniform ($\irate=1/\tauexp$, $L=\ln[\tauexp/(\tauexp-h)]$), scaled
($a/\tauexp$), power ($a\tauexp^{-b}$, $L=a[\tauexp^{1-b}-(\tauexp-h)^{1-b}]/(1-b)$), piecewise
constant over $\tauexp$ buckets with edges $2,4,\ldots,1024$\,h, and the power and piecewise forms
times 24 hour-of-day factors with mean one (``$\times$hour''). Each is fitted per category by
L-BFGS on up to 200{,}000 estimation steps. The BIC-best form is used downstream, and forms are
compared by holdout mean negative log-likelihood per step relative to uniform.

The standardized innovation
\begin{equation}\label{eq:dyn-u}
\ustd_t=\frac{\tilde\Zt_{t+h}-\tilde\Zt_t e^{L_t/2}}{\sqrt{e^{L_t}-1+k\,s_t^2/\varphi(\tilde\Zt_t)^2}}
\end{equation}
is $\N(0,1)$ under M0. We report its variance, excess kurtosis and tail frequencies.

The \emph{probit-drift} test regresses $\tilde\Zt_{t+h}-\tilde\Zt_t$ on $\tilde\Zt_t(e^{L_t/2}-1)$ for
steps with $L_t\le2$. M0 predicts a slope of one, the outward drift that keeps $\pr$ a martingale.
An IV version instruments with the previous hour's $\tilde\Zt$ to remove errors-in-variables bias
from quote noise.

The \emph{variance identity} is checked per market: $\sum(\Delta\pr)^2$ over two-sided quotes
(spread at most 4\cents) sampled every $k\in\{1,6,24\}$\,h, plus the final move
$(y-\pr_{\text{last}})^2$, compared with $\pr_0(1-\pr_0)$.

\subsubsection{Backbone, stochastic volatility, leverage and jumps}\label{sec:dyn-sv}

\paragraph{Backbone.}
For each backbone $g$ we fit by QLIKE
\begin{equation}\label{eq:dyn-backbone}
\E[(\Delta \pr)^2\mid \pr,\tauexp]=c_{j(\tauexp)}^2\,g(\pr)^2/g(1/2)^2+k\,s^2,
\end{equation}
with one scale per $\tauexp$ bucket $j$. The candidates are the M1 power family
$g=[\pr(1-\pr)]^\beta$, logit-normal ($\beta=1$), Jacobi ($\beta=1/2$) and the probit backbone
$\varphi(\Phi^{-1}(\pr))$ of M0; they are compared by holdout QLIKE relative to probit. A model-light
check regresses log cell means of $(\Delta\pr)^2$, Roll-corrected as
$\E(\Delta\pr_t)^2+2\E(\Delta\pr_t\Delta\pr_{t+1})$ \citep{roll1984}, on $\tauexp$-bucket fixed
effects and $2\beta\log[\pr(1-\pr)]$, with an event-level Poisson bootstrap (100 draws).

\paragraph{Volatility clustering.}
We estimate the within-market autocorrelation of $|\ustd_t|$ at lags 1--168\,h. It is computed
pooled and, removing static differences between markets, demeaned by market (markets with at least
48 steps). We fit $\Corr(|\ustd_t|,|\ustd_{t+k}|)=A\phi^k$ over $k\le24$ and report the half-life
$\ln0.5/\ln\phi$. Under log-normal stochastic volatility \citep{taylor1986},
$A=(x-1)/(\tfrac{\pi}{2}x-1)$ with $x=e^{\sigma_h^2/4}$ and $\sigma_h^2=\Var(\log\alpha^2)$. This
gives the M1 vol-of-vol $\nu=\tfrac12\sqrt{\sigma_h^2(1-\phi^2)}$ per $\sqrt{\text{hour}}$, which we
report per $\sqrt{\text{day}}$.

\paragraph{Leverage.}
The sign of $\rho$, in the $\sgn\Zt$ convention of M1, comes from
\begin{equation}\label{eq:dyn-lev}
\Corr\Big(\sgn(\Zt_t)\,\ustd_t,\;\tfrac16\textstyle\sum_{j=1}^{6}|\ustd_{t+j}|\Big),
\end{equation}
on steps with $\pr\in[0.1,0.9]$ and $\ustd$ winsorized at $\pm8$. A negative value means volatility
rises after moves toward $1/2$. A price-unit version replaces $\ustd$ by $\Delta\pr$ over the fitted
backbone volatility.

\paragraph{Jumps.}
Following M2, an EM \citep{dempster1977} fits the zero-mean scale mixture
\begin{equation}\label{eq:dyn-em}
\ustd\sim(1-\pi)\,\N(0,s_d^2)+\pi\,\N(0,s_d^2+s_J^2)
\end{equation}
to steps in which the mid moved; hours without a change are reported as the zero-move share. The
jump share of variance is $\sum\gamma_i(\Delta\pr_i)^2/\sum(\Delta\pr_i)^2$, with $\gamma_i$ the
posterior jump probability. The mixture is a random clock burst, which keeps $\pr$ an exact
martingale.

\subsubsection{Variance forecasts}\label{sec:dyn-forecast}

The target is $RV_{t,H}=\sum_{j<H}(\Delta\pr_{t+j})^2$ for $H\in\{1,6\}$\,h, on decision rows with at
least 24\,h of history. We compare five forecasts:
\begin{itemize}
\item \textsc{hist}: the trailing mean of $(\Delta\pr)^2$ over $W\in\{6,24,72\}$\,h, times $H$;
\item M0: the closed-form transition variance under the fitted clock;
\item \textsc{bb}: the backbone local variance of \cref{eq:dyn-backbone};
\item M1: the M0 variance times
  $\hat\alpha_t^2=(1-w)+w\,\mathrm{EWMA}_{\mathrm{hl}}[(\Delta\pr)^2/v_1]$ over the market's past
  steps, with $v_1$ the one-step M0 variance plus noise, $\mathrm{hl}\in\{3,12,48\}$\,h and
  $w\in\{0.25,\ldots,1\}$;
\item \textsc{m1bb}: \textsc{bb} times $\hat\alpha_t^2$.
\end{itemize}
Each forecast is calibrated on the \estim{} as $\kappa\cdot\text{base}+c_0s^2H+c_1$ by QLIKE, with
$W$, hl and $w$ chosen there. Holdout accuracy is measured by QLIKE, $\log f+RV/f$
\citep{patton2011}, and by MSE, with event-clustered $t$-statistics on per-row loss differences
\citep{dieboldmariano1995}.

\subsubsection{High-frequency in-game sample}\label{sec:dyn-hf}

For the nights of 7--8 and 8--9 November 2026 (17:00 to 06:00 UTC) we replay every recorded order-book
snapshot and delta of the public Kalshi websocket for game-winner series. Each market's book is
reduced to its top, persists across hourly files and is reset by snapshots. We keep one leg per
game, the first alphabetically excluding ties. Its in-game window runs from expected expiration
minus the league's game length to expected expiration plus 1\,h. The window is cut at the first
decided quote (mid at most 2\cents{} or at least 98\cents{}, spread at most 10\cents) and 10 minutes
before $T$. Legs with at least 30 minutes remaining are kept: 50 legs (NBA 14, NCAAF 14, NFL 10,
NHL 12), 6{,}965 in-game minutes and 267{,}422 top-of-book changes. % source: hf_results.json:games

The mid is sampled as of each point of an $h$-second grid, $h\in\{1,5,15,60,300\}$. Steps do not
overlap and use the uniform clock with $\tauexp$ in seconds. The start spread is at most 3\cents{}
(2\cents{} as a robustness check), the start mid is in $[0.03,0.97]$ and the end spread is at most
10\cents; $t$-statistics cluster by game.

\paragraph{Signature and jumps.}
The variance signature is $RV(h)/T$ per in-game hour for $h$ from 1 to 600\,s on the same windows
\citep{andersen2000signature}. Under i.i.d.\ noise of standard deviation $\eta$ in the mid,
\begin{equation}\label{eq:dyn-signature}
\E[RV(h)]/T=\sigma^2+2\eta^2/h ,
\end{equation}
which we fit across $h$ \citep{zhang2005,aitsahalia2005}.

At 5, 10 and 15\,s we compare realized variance with bipower variation,
$BV=\mu_1^{-2}\sum_i|r_i||r_{i-1}|$ with $\mu_1=\sqrt{2/\pi}$, and report the jump share
$(RV-BV)/RV$ \citep{barndorffnielsen2004,barndorffnielsen2006,huangtauchen2005}. Threshold jumps,
following \citet{leemykland2008}, have $|\ustd_i|>4\hat\sigma_i$ and $|\Delta\pr_i|\ge2$\cents,
where $\hat\sigma_i$ is the bipower scale of $\ustd$ over the preceding 30 minutes, floored at 0.25.

For a fixed score increment, Stern's Brownian game model predicts a jump size proportional to
$\varphi(\Zt)/\sqrt{\tauexp}$ \citep{stern1994,polsonstern2015}. We regress $\log|\Delta\pr|$ on
that quantity over the detected 10-second jumps, with league fixed effects; the predicted slope is
one. Leverage at 1, 2 and 5 minutes uses \cref{eq:dyn-lev} with the next five steps.

\subsubsection{Tick-grid control simulations}\label{sec:dyn-control}

\paragraph{Matched high-frequency control.}
For each leg we simulate an M0 path per second under the uniform clock from the leg's first mid. We
quote it on the 1\cents{} grid with the leg's observed spread, bid
$=\operatorname{round}(100\pr-s/2)$\cents{} and ask $=$ bid $+\,s$, apply the same cuts, and compute
the same statistics. The control therefore has M0, ticks and the real spreads, but no jumps and no
stochastic volatility.

\paragraph{Hourly drift control.}
M0 paths are simulated on each market's real hourly grid with the category's fitted clock, in three
versions:
\begin{itemize}
\item \emph{true}: not quoted;
\item \emph{tick}: quoted on the grid with the real spreads;
\item \emph{stale}: quoted on the grid but re-quoted only in hours with a real candle.
\end{itemize}
To locate any shortfall we refit the drift regression inside $|\ustd|\le2$, a window centred on M0's
mean that biases the slope toward one, and inside $|\Delta\Zt|\le2$\,sd, centred on zero, which
biases it toward zero. On $|\Zt|\ge0.5$ we also compute the local second-order martingale ratio
$\E[\sgn(\Zt)\Delta\Zt]/\E[\tfrac12|\Zt|(\Delta\Zt)^2]$, which is one when $\pr$ is a martingale,
whatever the clock.
